% =====================================================================
%  Capítulo 6: Conclusiones
%  Rúbrica: 5 %. Conexión explícita con cada objetivo; limitaciones;
%  implicancias prácticas; trabajo futuro concreto.
%  Las conclusiones se apoyan en los RESULTADOS PROPIOS (no en citas).
%  Reescrito 2026-11-20 según los objetivos aprobados (VF-V2).
% =====================================================================
\chapter{Conclusiones}
\label{cap:conclusiones}

\section{Conclusiones}
\label{sec:conclusiones}

Se cumplió el objetivo general de diseñar un sistema experimental de humedales construidos de flujo vertical para tres fuentes de aguas residuales en Frutícola SABA SpA, integrando su cálculo hidráulico y estructural.
El diseño se compone de tres líneas independientes de \qty{6.00}{\square\meter} de área activa en módulos IBC, alimentadas por pulsos desde estanques dosificadores ubicados sobre una plataforma de acero, y queda documentado en memorias de cálculo, planos, especificaciones técnicas e itemizado.
\revisar{Confirmar que el objetivo se cumple completo cuando resuelva las preguntas pendientes con el profesor.}

En relación con el primer objetivo específico, la caracterización mostró que las tres fuentes son distintas y deben tratarse por separado.
Las aguas servidas presentan una DQO de 338 a \qty{411}{\milli\gram\per\liter}, propia de las aguas servidas rurales, pero un nitrógeno anormalmente bajo que debe confirmarse; las aguas de lavado de agroquímicos son variables (DQO de 76 a \qty{393}{\milli\gram\per\liter}) y de producción limitada; y el agua del estero tiene una carga orgánica mínima, pero coliformes fecales que la dejan fuera de los usos de riego sin tratamiento.
El caudal de diseño de la línea doméstica es de \qty{0.75}{\cubic\meter\per\day} (\qty{0.60}{\cubic\meter\per\day} de la casa y \qty{0.15}{\cubic\meter\per\day} de las oficinas), mientras que los caudales del estero y de los agroquímicos quedan por definir con un criterio distinto al de igualar la carga orgánica.

En relación con el segundo objetivo específico, el balance de oxígeno verificó que \qty{6.00}{\square\meter} por línea aportan casi tres veces el oxígeno que demanda la línea doméstica, con una carga de \qty{30}{\gram\per\square\meter\per\day} de DQO y \qty{125}{\milli\meter\per\day}, superiores a las de las guías europeas, pero dentro de la experiencia en climas cálidos.
En el cálculo hidráulico, la línea del estero es la crítica: con tubería de \qty{25}{\milli\meter} requiere una altura de unos \qty{31}{\meter}, fuera del alcance de las bombas consideradas, por lo que su diámetro debe aumentarse a \qty{40}{\milli\meter} o rediseñarse el llenado.
El bajo caudal del microrriego exige un circuito de retorno en la impulsión.

En relación con el tercer objetivo específico, los perfiles cajón de la plataforma resisten las cargas de operación con amplio margen, y el diseño queda controlado por la deformación y por la estabilidad de las fundaciones.
El tablero de \qty{15}{\milli\meter} debe aumentarse a \qty{21}{\milli\meter} o apoyarse en vigas secundarias más próximas para conservar la horizontalidad que necesitan los sifones.
Las zapatas de $\num{1.00} \times \num{1.00} \times \num{0.50}$~\unit{\meter} cumplen con holgura frente a las cargas estáticas, pero su verificación sísmica solo es satisfactoria si los seis pilares comparten el corte sísmico.
En consecuencia, la horizontalidad de la plataforma y la estabilidad de las zapatas frente al sismo son las condiciones críticas del diseño.

\section{Limitaciones}
\label{sec:limitaciones}

El diseño se basa en seis campañas de muestreo realizadas en enero y marzo de 2026, con solo dos en la casa, por lo que no representa la variación estacional de las fuentes; además, el nitrógeno medido en las aguas servidas es incoherente con la literatura, y el caudal de las oficinas se apoya en un mes de registro.
Los caudales del estero y de las aguas de lavado no pudieron definirse con la información disponible, y el volumen de lavado que genera el predio no ha sido cuantificado.

En lo estructural, no se dispone de un estudio de mecánica de suelos, por lo que la capacidad admisible y el ángulo de fricción del suelo son supuestos, y la verificación sísmica depende de parámetros aún no confirmados (versión de la NCh2369, categoría de ocupación, factor de modificación de la respuesta y comportamiento del tablero como diafragma).
Finalmente, el diseño no ha sido validado en operación: la eficiencia esperada del humedal, en particular la remoción de coliformes en la línea del estero, es una estimación basada en la literatura.

\section{Trabajo futuro}
\label{sec:trabajo-futuro}

En el corto plazo, se recomienda completar el diseño antes de la construcción: ampliar la caracterización a otras estaciones del año y verificar las unidades del nitrógeno, medir el volumen real de lavado de equipos del predio, definir los caudales de las líneas del estero y de agroquímicos, y realizar un estudio de mecánica de suelos para recalcular las zapatas con la NCh2369:2025.

Una vez construida la estación, su operación permitirá contrastar el diseño con datos reales: medir la eficiencia de cada línea, calibrar el balance de oxígeno, evaluar la carga de la Fase~2 (\qty{60}{\gram\per\square\meter\per\day} de DQO) sin colmatación y comprobar si los efluentes cumplen los requisitos de calidad para riego, en coordinación con los estudios agronómicos del proyecto ANID/Anillos/ATE250033 sobre el cultivo del cerezo.
